\clearpage
\VerdierPhysicalPageAnchor{0235}%
\begingroup\small\setlength{\parskip}{0pt}\setlength{\emergencystretch}{1.5em}
\setlength{\abovedisplayskip}{4pt plus 2pt minus 2pt}%
\setlength{\belowdisplayskip}{4pt plus 2pt minus 2pt}%

\VerdierSemanticAnchor{ch4d-4-2-2-cont}
\noindent ここで $(p,q)$ と $(p',q')$ は
$\mathcal{F}\ell_1(\widetilde{\mathbb Z})$ の元で、$p\leq p'$ かつ
$q\leq q'$ を満たす。これは $\operatorname{comp}(\mathcal A)$ の標準射

\VerdierSemanticAnchor{ch4d-4-2-2-3}
\begin{equation}
\psi\bigl((p,q),(p',q')\bigr)\colon
X(p,q)\longrightarrow X(p',q').
\tag{4.2.2.3}
\end{equation}
の $D(\mathcal A)$ における像である。こうして
$\mathcal{F}\ell_1(\widetilde{\mathbb Z})$ から $D(\mathcal A)$ への関手が
定まった。$\mathcal{F}\ell_2(\widetilde{\mathbb Z})$ の各元 $(p,q,r)$ に対し、

\VerdierSemanticAnchor{ch4d-4-2-2-4}
\begin{equation}
S(p,q,r)\colon
0\longrightarrow X(p,q)\longrightarrow X(p,r)
\longrightarrow X(q,r)\longrightarrow0
\tag{4.2.2.4}
\end{equation}
を、射が $X$ のフィルトレーションから得られる標準射である
$\operatorname{comp}(\mathcal A)$ の完全列とする。最後に、
$\mathcal{F}\ell_1(\widetilde{\mathbb Z})$ の各元 $(p,q)$ に対し、

\VerdierSemanticAnchor{ch4d-4-2-2-5}
\begin{equation}
J(p,q)\colon\mathcal X(p,q)\longrightarrow X(p,q)
\tag{4.2.2.5}
\end{equation}
を、$\operatorname{comp}(\mathcal A)$ の射

\VerdierSemanticAnchor{ch4d-4-2-2-6}
\begin{equation}
j(p,q)\colon\mathcal X(p,q)\longrightarrow X(p,q),
\tag{4.2.2.6}
\end{equation}
の $D(\mathcal A)$ における像とする。その成分は次である：

\VerdierSemanticAnchor{ch4d-4-2-2-7}
\begin{equation}
\left\{
\begin{aligned}
j(p,q)^i\colon\mathcal X(p,q)^i
 &=X(q)^i\oplus X(p)^{i+1}
   \longrightarrow X(q)^i/X(p)^i=X(p,q)^i,\\
j(p,q)^i&=(\pi^i,0),\\
&\text{ここで $\pi^i$ は標準エピ射である：}\\
\pi^i\colon X(q)^i&\longrightarrow X(q)^i/X(p)^i.
\end{aligned}
\right.
\tag{4.2.2.7}
\end{equation}

\VerdierSemanticAnchor{ch4d-4-2-3}
\medskip\noindent\textbf{命題 4.2.3.}\quad
\textit{$\mathcal{F}\ell_2(\widetilde{\mathbb Z})$ の各元 $(p,q,r)$ に対し、}
\[
\delta(p,q,r)\colon X(q,r)\longrightarrow X(p,q)[1]
\]
\noindent\textit{を $D(\mathcal A)$ の射}

\VerdierSemanticAnchor{ch4d-4-2-3-1}
\begin{equation}
\delta(p,q,r)=\delta\bigl(S(p,q,r)\bigr),
\tag{4.2.3.1}
\end{equation}
\noindent\textit{とする。ここで $\delta\bigl(S(p,q,r)\bigr)$ は、完全列 $S(p,q,r)$
(4.2.2.4) に標準 $\delta$-関手
$\operatorname{comp}(\mathcal A)\longrightarrow D(\mathcal A)$ (1.3.2)
が対応させる $D(\mathcal A)$ の射である。}

\VerdierSemanticAnchor{ch4d-4-2-3-ab}
\noindent\textit{a) 対象 $X(p,q)$ (4.2.2.1)、射
$\varphi\bigl((p,q),(p',q')\bigr)$ (4.2.2.2)、および射
$\delta(p,q,r)$ (4.2.3.1) は、$D(\mathcal A)$ に値をとる
スペクトル対象を定める。}

\noindent\textit{b) 射 $J(p,q)$ (4.2.2.5) は、スペクトル対象
$\bigl(\mathcal X(p,q),\delta(p,q,r)\bigr)$ からスペクトル対象
$\bigl(X(p,q),\delta(p,q,r)\bigr)$ への同型を定める。}

\VerdierSemanticAnchor{ch4d-4-2-3-proof}
a) は標準 $\delta$-関手
$\operatorname{comp}(\mathcal A)\longrightarrow D(\mathcal A)$ の性質
(1.3.1) から直ちに従う。b) を証明するには (1.3.2) を参照されたい。
\endgroup
